\chapter{WCCI36: Target Environment and Controls}
\label{ch:wcci-transfer}

The experiments now move from \texttt{bus14} to the larger WCCI 2020 grid.
This chapter defines the target setting and checks which models can learn on it
before testing transfer.

\section{Target environment}
\label{sec:wcci-environment}

WCCI has more substations, lines and agents than \texttt{bus14}
(Table~\ref{tab:wcci-env-comparison}). The agents still control separate areas,
but their areas are not equal in size.

\begin{table}[H]
  \centering
  \small
  \caption{Source and target environments. Both use topology actions,
  decentralized agents and the same train/test split rule.}
  \label{tab:wcci-env-comparison}
  \setlength{\tabcolsep}{6pt}
  \begin{tabular}{L{0.30\textwidth}L{0.28\textwidth}L{0.30\textwidth}}
    \toprule
    & \texttt{bus14} & \texttt{bus36\_wcci\_nomaint} \\
    \midrule
    Substations & 14 & 36 \\
    Lines & 20 & 59 \\
    Agents & 3 & 4 \\
    Chronics & 1{,}004 & 2{,}880 \\
    Episode length & up to 8{,}064 steps & up to 8{,}062 steps \\
    Maintenance & none & disabled (Section~\ref{sec:wcci-nomaint}) \\
    \bottomrule
  \end{tabular}
\end{table}

\subsection{Removing maintenance}
\label{sec:wcci-nomaint}

The WCCI environment normally includes scheduled line maintenance, while
\texttt{bus14} does not. Maintenance is disabled here so that a transferred
policy is not tested on a hazard absent from its source environment. This also
keeps the observation spaces aligned.
\\
\\
On the full WCCI test split, removing maintenance raises do-nothing mean
survival from $27.3\%$ to $57.3\%$ (Table~\ref{tab:wcci-maintenance-effect}).
The change is large enough that maintenance-on and maintenance-off results
should not be compared directly.

\begin{table}[H]
  \centering
  \small
  \caption{Do-nothing performance on the full WCCI test split. The 576
  chronics are the same in both rows.}
  \label{tab:wcci-maintenance-effect}
  \setlength{\tabcolsep}{7pt}
  \begin{tabular}{lrrr}
    \toprule
    Setting & Mean survival & Median survival & Complete episodes \\
    \midrule
    Maintenance enabled  & 27.3\% & 12.1\%  & 39/576 (6.8\%) \\
    Maintenance disabled & 57.3\% & 100.0\% & 303/576 (52.6\%) \\
    \bottomrule
  \end{tabular}
\end{table}

\subsection{Reduced action spaces}
\label{sec:wcci-action-space}

The full decentralized action space contains almost 67{,}000 actions and is
strongly unbalanced across agents. The candidate-action models are therefore
evaluated with reduced spaces of $k\in\{32,64,128,256\}$. Here, $k$ is a cap:
agents with fewer valid actions keep their complete list.
\\
\\
The greedy survival values below are not computed on the full
576-chronic WCCI test split. They use only the 24 difficult chronics from the
fixed 50-chronic evaluation subset defined in
Section~\ref{sec:wcci-evaluation-protocol}.

\begin{table}[H]
  \centering
  \small
  \caption{Per-agent action counts on \texttt{bus36\_wcci\_nomaint} and
  one-step greedy mean survival on a subset of the WCCI test split:
  the 24 difficult chronics from the fixed 50-chronic evaluation set. These are
  not full-test averages. Agent headings give the number of controlled
  substations.}
  \label{tab:wcci-action-space}
  \setlength{\tabcolsep}{5pt}
  \begin{tabular}{lrrrrrr}
    \toprule
    Action space
      & \texttt{agent\_0} (5)
      & \texttt{agent\_1} (5)
      & \texttt{agent\_2} (9)
      & \texttt{agent\_3} (17)
      & Total
      & Greedy mean survival \\
    \midrule
    Full list & 77 & 65{,}642 & 127 & 1{,}119 & 66{,}965 & -- \\
    $k=32$   & 32 & 32  & 32  & 32  & 128 & 33.1\% \\
    $k=64$   & 64 & 64  & 64  & 64  & 256 & 67.9\% \\
    $k=128$  & 77 & 128 & 127 & 128 & 460 & 68.2\% \\
    $k=256$  & 77 & 256 & 127 & 256 & 716 & 94.9\% \\
    \bottomrule
  \end{tabular}
\end{table}

\noindent A one-step greedy policy is used only as a reference for what each reduced
space contains. Its difficult-cohort mean survival rises from $33.1\%$ at $k=32$ to
$94.9\%$ at $k=256$. A larger action list therefore gives more useful choices,
but it also gives the learned policy a harder selection problem. This trade-off
is measured in Chapter~\ref{ch:screen-f-transfer}.

\section{Evaluation protocol}
\label{sec:wcci-evaluation-protocol}

All reported transfer evaluations use the same 50 WCCI chronics. Do nothing
completes 26 of them, giving a 52\% completion rate. These 26 form the \emph{easy} cohort; the remaining 24 form
the \emph{difficult} cohort. Do-nothing mean survival is $56.0\%$ overall and
$8.34\%$ on the difficult cohort.
\\
\\
Reporting both values separates two goals: preserving cases that already need
no action and improving cases where inaction fails. The number of difficult
chronics completed is reported as \emph{rescues}. The one-step greedy policy is
a strong reference, not an upper bound, because it queries the simulator at
every step.

\section{Controls trained on WCCI}
\label{sec:wcci-architecture-choice}

Figure~\ref{fig:wcci-target-controls} compares three controls trained directly
on WCCI: the flat-MLP baseline, the four selected GNN encoders with
grid-specific MLP action heads, and the eight GNN candidate-scorer variants.
All architecture cells use seed 0. For a family with several cells, the bar
reports their arithmetic mean and the whisker spans their minimum and maximum. The
number above each bar gives the number of cells included. Figure~\ref{fig:wcci-target-controls}
keeps only \(k=64\), because it is roughly the same size the action space on the bus 14 environment.

\noindent\textcolor{green!50!black}{\textbf{Incomplete settings excluded:}
at \(k=32\), only the flat MLP is currently available; at \(k=256\), the flat
MLP and selected GNN encoder with an MLP head are available.}

\begin{figure}[H]
  \centering
  \includegraphics[width=\linewidth]{figures/wcci_target_controls.png}
  \caption{Matched comparison of WCCI controls trained from scratch at
  \(k=64\). Bar colour identifies the model family. Bar height is the family
  mean, whiskers show the minimum--maximum range when several architecture
  cells are available, and labels give the number of cells. The dashed line is
  do nothing.}
  \label{fig:wcci-target-controls}
\end{figure}

\noindent At the matched \(k=64\) setting, the overall mean across architecture
cells is \(32.64\%\) for the flat MLP, \(38.40\%\) for the four fixed-head GNNs and \(47.21\%\) for
the eight candidate scorers. All three remain below the \(56.0\%\) do-nothing
reference. These values are surprisingly low for models
retrained directly on the target grid. On the difficult cohort, however, the candidate scorer has the
strongest mean at \(10.51\%\), compared with \(8.82\%\) for the flat MLP and
\(5.26\%\) for the fixed-head GNNs. Its best cell reaches \(18.46\%\).
\\
\\
The candidate scorer is also the only graph actor here whose complete policy
can cross grids. Its shared scorer is applied to each candidate, so its output
size does not depend on the number of agents or actions. A fixed-list head has
one output per action index and cannot be reused unchanged on WCCI.
\\
\\
Exact flat-MLP and fixed-list GNN results are given in
Appendix~\ref{app:wcci-result-atlas}.

\begin{tcolorbox}[
  colframe=red!75!black,
  colback=red!5!white,
  title=Chapter 7 Summary: Target Environment and Controls]
\textbf{Purpose:} Define the WCCI target and scratch-trained
reference controls for the transfer study.
\\
\\
\textbf{Content:} We quantify the maintenance effect, reduce
the action space, define a fixed 50-chronic evaluation, and compare the three
control families at the matched $k=64$ setting.
\\
\\
\textbf{Findings:}
\begin{itemize}
  \item Disabling maintenance raises do-nothing mean survival from $27.3\%$ to
  $57.3\%$. On the 24 difficult cases, it reaches only $8.34\%$.
  \item Greedy difficult-case survival rises from $33.1\%$ at $k=32$ to
  $94.9\%$ at $k=256$, but larger lists make ranking harder.
  \item Direct retraining is surprisingly weak: all three
  $k=64$ family means remain below do nothing overall. The candidate scorer is
  strongest on difficult cases, with a $10.51\%$ mean and an $18.46\%$ best
  cell.
\end{itemize}
\textbf{Takeaway:} Target-grid retraining alone does not solve
WCCI control. These controls provide the reference for encoder reuse and
complete-policy transfer.
\end{tcolorbox}
